
\subsection{An actual finite colimit of cut-surface groupoids}

\begin{result}{B.27}{Proved}{A strict finite colimit with the five elementary types}
Let $S$ be an oriented stable surface with individually labelled
boundary circles $I$, let $H=\Gamma(S,\partial S)$, and put
$A=\mathbb Z^I$, with its positive boundary-twist generators.
There is a strict finite diagram of finite-object groupoids,
with a free action of $\mu=A/2A$ on their object sets, whose
ordinary colimit is equivalent to the fundamental groupoid of
the full tangent-ray bundle of the corresponding moduli stack,
restricted to the specified finite rays.  The equivalence
preserves the boundary half-turns and full turns.

The pieces are obtained by cutting and sewing from the five
types
\[
 (0,3),\quad(0,4),\quad(0,5),\quad(1,1),\quad(1,2). \tag{B.74}
\]
For genus zero the index can be the finite poset of reduced
labelled trees $G$ with
\[
 \sum_v(\operatorname{val}(v)-3)\leq2,
\]
ordered by edge contraction.  The local five-point and elliptic
pieces here are the explicitly presented groupoids of B.19,
B.25 and B.9--B.11.  For a pinned graph with several vertices,
finite products and the flag-to-edge induction of B.1 give the
same result with the full external lattice
$\mathbb Z^{E(G)\amalg I(G)}$.

This proves a geometric finite-colimit recipe with all five
pieces (B.74).  It does not remove the elliptic pieces in favour
of only three- and four-point sphere pieces, nor identify an
unstated original finite ray convention.  Those are the stronger
comparisons retained in B.8 and B.15.
\end{result}
\begin{proof}
We give the index, all functors, the strictification and the
comparison inverse.  In particular, the colimit here
means an ordinary colimit of a specified strict diagram.

Put $d=3g-3+|I|$ and let $X=|P_2|$, with $P_2$ the
multicurves having at least $d-2$ components.  When $d\leq2$
the empty multicurve is allowed.  The simple connectedness
theorem proved and audited in D.5 applies to this precise
$P_2$.  For $d\leq2$ the asserted construction is one of
(B.74), already explicitly computed.  For $d\geq3$, every
vertex of $P_2$ is a nonempty multicurve.  Cutting a multicurve
$\sigma$ gives
\[
 \sum_Q(3g(Q)-3+n(Q))=d-|\sigma|\leq2.             \tag{B.75}
\]
Each term is nonnegative.  The stable types of nonnegative
complexity at most two are exactly (B.74).

Use the barycentric action category: its cells are flags of
multicurves in $P_2$, and its arrows are face inclusions.
Take one representative of every $H$-orbit of cells and every
incidence orbit.  The resulting index $\mathcal I$ is finite:
the dual stable graph has bounded numbers of vertices and
edges, and a flag is a chain of at most three edge subsets.
If a face has several incidences they remain separate arrows.
There is no quotient of those arrows to connected components.
Cell stabilizers are the stabilizers of the full flags, so
they fix the barycentric cells pointwise.

For a representative cell $i$, write $K_i\subset H$ for
its actual stabilizer.  It contains $A$ centrally.  Cutting
all curves of its finest multicurve describes $K_i$ using
the products of the component boundary-fixed groups, the
identifications of paired boundary twists, and the finite
automorphisms preserving the flag.  The paired twist map
is exactly the map of B.1, not a newly postulated quotient.
All component groups have the explicit presentations (B.74).
For finite graph permutations retain their actual lifts and
multiplication cocycles.  One can construct their finite
presentation by choosing one lift of each graph permutation
and adding its conjugation and multiplication relations in
the cut groups.  These relations are computed by restriction
to the labelled components and collars; the component word
problems are the braid and free-group word problems already
specified in B.9, B.18 and B.19.  This procedure makes no
splitting assertion for a finite graph-automorphism extension.
The nontrivial lift and flag calculations of B.11 and D.27
are explicit instances of this construction.

Choose actual incidence transporters $g_\alpha\in H$.
They send the chosen lifted face to its representative.
For an arrow $\alpha:i\to j$ the associated local functor
on vertex groups is
\[
 \phi_\alpha:K_i\longrightarrow K_j,\qquad
 k\longmapsto g_\alpha k g_\alpha^{-1}.            \tag{B.76}
\]
When an incidence stabilizer is required, $i$ denotes
that full flag, so (B.76) is defined on precisely that
stabilizer.  The transporter choices for a composition
give the natural isomorphism
\[
 c_{\beta,\alpha}
   =g_{\beta\alpha}(g_\beta g_\alpha)^{-1}:
 \phi_\beta\phi_\alpha\Longrightarrow\phi_{\beta\alpha}.
                                                               \tag{B.77}
\]
The element is in the target stabilizer.  Substitution
shows its naturality and its associativity identity:
both ways of comparing three composed conjugations
cancel to $g_{\gamma\beta\alpha}
(g_\gamma g_\beta g_\alpha)^{-1}$.
For identity arrows choose transporter $1$.
Thus this is an actual unital pseudofunctor, including
its specified natural transformations.

Here is its framed version.  Write $E=\tfrac12A$ and
identify $\mu=E/A$ by the indicated basis.  Define
\[
 K_i^{1/2}=(K_i\times E)/
       \{(a,-a):a\in A\},\qquad
 \chi_i:K_i^{1/2}\longrightarrow E/A.             \tag{B.78}
\]
The original $K_i$ injects and is the kernel of $\chi_i$.
Let $F(i)$ be the translation groupoid on the finite
set $\mu$: an element $h\in K_i^{1/2}$ is an arrow
$\epsilon\to\epsilon+\chi_i(h)$.  Equations (B.76)--(B.77)
extend by the identity on $E$ and hence give equivariant
functors and natural transformations on these groupoids.
The positive element $\frac12e_\ell$ gives an arrow
between opposite leg rays, and its square is exactly
$e_\ell=z_\ell$.  Translation on $\mu$ is free.

This formal description is the concrete tangent-ray
groupoid: choose a base framing and rotate each boundary
coordinate through zero or a positive half-turn.  A
path between two such framings is an element of (B.78),
and composition is concatenation.  B.1 identifies the
same construction under sewing.  If local finite models
retain additional node-ray objects, choose the positive
half-turn paths from one node-ray choice to the others
and use
\[
 f\longmapsto \ell_{t(f)}^{-1}f\,\ell_{s(f)}.       \tag{B.79}
\]
This is an equivalence to the stated stabilizer model.
For one node its first half-turn maps to $1$ and its
second to the full node twist.  Thus (B.79) retains
the meridian and does not quotient its lattice.
It is equivariant for all external leg rays.

We now replace the pseudofunctor by a strict one without
losing finiteness or freeness.  Define $\widetilde F(i)$
to have objects
\[
 (j,\alpha:j\to i,x),\qquad x\in\operatorname{Ob}F(j),
\]
and morphisms
\[
 \operatorname{Hom}_{\widetilde F(i)}
 ((j,\alpha,x),(k,\beta,y))
 =\operatorname{Hom}_{F(i)}(F(\alpha)x,F(\beta)y).
                                                               \tag{B.80}
\]
Composition is that of $F(i)$.  For $u:i\to l$, the
object map appends $u$, and the arrow map is
\[
 f\longmapsto c_{u,\beta}\,F(u)(f)\,c_{u,\alpha}^{-1}.
                                                               \tag{B.81}
\]
The associativity identity checked after (B.77) proves
that two successive maps (B.81) equal the map for their
composite, on the nose.  Thus $\widetilde F$ is a strict
finite diagram.  Evaluation into $F(i)$ is fully faithful
and essentially surjective, with inverse on objects
$x\mapsto(i,1_i,x)$ and the identity comparison arrows
from (B.80).

There is also a direct description of its ordinary colimit.
Its objects are the pairs $(i,x)$; the object
$(j,\alpha,x)$ is the image of $(j,1_j,x)$.
For each incidence $\alpha:j\to i$, the identity morphism
of $F(\alpha)x$ in (B.80) supplies an isomorphism from
$(j,x)$ to $(i,F(\alpha)x)$.  The remaining relations
are precisely the local groupoid relations, naturality
of these isomorphisms, and (B.77).  Conversely these
generators define a cone on (B.80)--(B.81), by the
displayed identity comparison arrows.  The two universal
maps fix every local morphism and every incidence
isomorphism, and are inverse.  This explicitly identifies
the ordinary colimit with the groupoid presentation of
the specified pseudofunctor.

Freeness is now checked for this particular colimit.
Every object map (B.81) preserves the tags $j,x$.
Consequently the object identifications never identify
two different elements of the sign torsor attached to
one $j$.  The colimit has the free sign action on its
objects $(j,\epsilon)$.  This is the additional argument
missing from the unrestricted colimit inference refuted
in B.5.

To verify the comparison and its inverse, let $H^{1/2}$
be (B.78) with $K_i$ replaced by $H$.  It acts on $X$
through $H/A$, so the half-boundary turns act trivially
on $X$.  Form the action-path groupoid whose objects are
$(v,\epsilon)$, with $v$ a vertex of the barycentric
subdivision of $X$, and whose arrows are
\[
 (h,[p]):(v,\epsilon)\longrightarrow(w,\epsilon+\chi(h)),
 \quad p\text{ a path from }hv\text{ to }w.
                                                               \tag{B.82}
\]
Use endpoint-preserving homotopy classes of paths.
Composition is $(k,[q])(h,[p])=(kh,[q*(kp)])$.
The local groupoids and incidence isomorphisms above
map to (B.82) by their actual transporter and cell paths.
A path in the cell complex reads conversely as a word
in those generators.  Backtracking cancels a generator
with its inverse; crossing a two-simplex gives exactly
(B.77).  This proves that the two word maps are inverse
on the full subgroupoid of the chosen orbit representatives.
It is the groupoid form of the action-complex proof,
with its generator and relation maps retained.

The functor from (B.82) to the translation groupoid of
$H^{1/2}$ sends $(h,[p])$ to $h$.  It is fully faithful:
for every $h$ the required path exists because $X$
is connected, and any two choices are homotopic because
$X$ is simply connected.  An inverse sends every sign
object to $(v_0,\epsilon)$ for a chosen reference vertex
$v_0$, and sends $h$ to $(h,[p_h])$, where $p_h$ is
any edge path from $hv_0$ to $v_0$.  The preceding
word-reading procedure computes its representative in
the finite colimit.  The choice has no effect on that
morphism, and concatenation proves the multiplication
law.  Chosen paths from the other representative
vertices to $v_0$ give the natural isomorphism from
the other composite to the identity.  In local
coordinates the forward functor sends a morphism
$f$ of (B.80) to
\[
 g_\beta^{-1}\,f\,g_\alpha
\]
in $H^{1/2}$.  This is the announced actual comparison,
including its inverse and all conjugating transports.

In genus zero, an individually leaf-labelled reduced
tree has no nonidentity automorphism: each internal
edge is determined by its two leaf sets, and the
vertices are then determined by incidence.  A graph
contraction is determined by the subset of edges retained.
Any two flags with the same labelled trees are related
by a surface homeomorphism preserving those edge
partitions.  Hence the orbit flags are exactly the
chains in the finite tree poset stated in the theorem.
Moreover the stabilizer of a finer tree fixes every
edge separately, so it preserves every contracted face.
We may therefore use that poset itself with (B.76)--(B.77);
its barycentric action presentation is the same one
just proved.  Strictification (B.80) does not change
its index.

Finally, for a pinned graph with several vertices,
apply the construction at each vertex, take the finite
product diagram, and apply the central induction and
ray-torsor tensor map of B.1 to each value and arrow.
Products of groupoids preserve colimits separately,
as follows from their adjunction with functor groupoids.
Central induction preserves them by its defining
extension-of-operators adjunction; the explicit
flag-to-edge map adds the two exponents at a sewn node.
Thus the resulting colimit is the actual sewn groupoid,
and its comparison is the sewn product of the already
specified functors.  Reapplying the sign-object
construction with the external lattice
$\mathbb Z^{E(G)\amalg I(G)}$ gives the same freeness
proof and the required node and leg half-turns.
Finite graph symmetries can be retained as an action
groupoid on these pinned objects; this introduces
arrows and never takes their coarse object quotient.
\end{proof}
